\documentclass[10pt,a4paper]{amsart}
\usepackage[margin=19mm]{geometry}
\usepackage{amsmath,amssymb,mathtools}
\usepackage[scr=rsfs,cal=euler]{mathalfa}
\usepackage{xfrac}
\usepackage[unicode,hidelinks,bookmarks=false]{hyperref}
\newcommand{\ie}{i.e.\ }
\newcommand{\cf}{cf.\ }
\newcommand{\resp}{respectively\ }
\newcommand{\Subsection}{\subsection}
\providecommand{\nobreakdash}{\nobreak\hbox{-}}
\setlength{\emergencystretch}{2em}
\multlinegap0pt
\usepackage{amssymb}
\usepackage{tensor}
\usepackage{mathtools}
\usepackage{dsfont}
\usepackage[shortlabels]{enumitem}
\newtheorem{lem}{Lemma}
	\newcommand{\DeltaH}{\Delta_{\H}}
\renewcommand{\H}{\mathrm{H}}
\newcommand {\T}	{\mathbb{T}}
\DeclareMathOperator{\divH}{div_{\H}}
	\newcommand{\nablaH}{\nabla_{\H}}
	\newcommand{\rH}{\H}
	\newcommand{\rL}{\mathrm{L}}
\title{The Ekman spiral in primitive equations with wind-driven boundary conditions and Coriolis force}
\author{Tim Binz}
\date{Source: arXiv:2507.20031v2 (CC BY 4.0)}
\begin{document}
\maketitle

\noindent\emph{Source and licence.} The Ekman spiral in primitive equations with wind-driven boundary conditions and Coriolis force. Version arXiv:2507.20031v2 (CC BY 4.0), distributed by arXiv under the Creative Commons Attribution 4.0 International licence, http://creativecommons.org/licenses/by/4.0/, as recorded in the arXiv licence metadata for this exact version and shown on the abstract page https://arxiv.org/abs/2507.20031v2. Full licence notice: \texttt{licenses/arxiv-2507.20031-CC-BY-4.0.txt}.\par\smallskip

\noindent\textit{Editorial scope.} Counted unit: the article's lem NSE\_2 (source0.tex lines 2238--2255). The article's complete original proof of it is delivered with it (source0.tex lines 2256--2311, one \texttt{proof} environment of 56 lines). No mathematics is written, repaired or re-derived: the statement and the proof are the article's own lines, in the article's own notation, and no retained line is substituted or re-flowed.  No further source-local context is needed: every cross-reference of every retained line resolves inside the two delivered spans.  Nothing was omitted from any retained span, so the package carries no omission record.  No retained line contains a citation, so the wrapper carries no bibliography and no bibliography entry.  No retained line contains a figure environment, an image-inclusion command or drawing code, so no image is delivered and the package declares no asset dependency.  Editorial formatting only, in addition: an AMS \texttt{amsart} preamble that restates the article's own declarations and macros used by the retained lines, namely the environments \texttt{lem} on separate counters, together with the standard packages those lines need.  The retained environments print in this document as Lemma 1 in order of appearance, whereas the article prints the same items with the numbering of its own full document.  The complete native source of the article as distributed in the arXiv e-print for this exact version is bundled unchanged as \texttt{sources/182258/source0.tex}.
\begin{lem}[Estimates of the 2D Stokes equations]\label{lem:2D NSE_2}
	Let $(u,\pi)$ be the solution of the two-dimensional Stokes equations with rotation
	\begin{equation*} 
		\left\{
		\begin{aligned}
			\partial_t u - \nu_{\rH} \DeltaH u + \nablaH \pi + u^\perp &= k,\\
			\divH u &= 0
		\end{aligned}
		\right. 
	\end{equation*} 
	for $k \in \rL^2(0,T;\rL^2(\T^2))$. Then $(u,\pi)$ satisfies the estimate 
	\begin{equation*}
		\sqrt{\nu_{\rH}} \cdot \partial_t \| \nablaH u \|_{\rL^2(\T^2)}^2
		+ \nu_{\rH} \cdot \| \DeltaH u \|_{\rL^2(\T^2)}^2
		+ \| \nablaH \pi \|_{\rL^2(\T^2)}^2
		\leq C \cdot (\| h \|_{\rL^2(\T^2)}^2+\| u \|_{\rL^2(\T^2)}^2) .
	\end{equation*}
\end{lem}

\begin{proof}
	Multiplying by $\partial_t u - \DeltaH u$, integration by parts and Young's inequality yield
	\begin{align*}
		\| \partial_t u \|_{\rL^2(\T^2)}^2 
		+
		\sqrt{\nu_{\rH}}\partial_t \| \nablaH u \|_{\rL^2(\T^2)}^2
		+\nu_{\rH} \|\DeltaH u \|_{\rL^2(\T^2)}^2
		\leq \int_{\T^2} k (\partial_t u - \DeltaH u) \\
		\leq 
		C \| k \|_{\rL^2(\T^2)}^2
		+
		\varepsilon \| \partial_t u \|_{\rL^2(\T^2)}^2 
		+
		\varepsilon 
		\|\DeltaH u \|_{\rL^2(\T^2)}^2
		, 
	\end{align*}
	where we have used 
	\begin{align*}
		-\int_{\T^2} \nablaH \pi \DeltaH u
		&=  \int_{\T^2} \pi \divH \DeltaH u
		= \int_{\T^2} \pi \DeltaH \divH u
		= 0, \\
		-\int_{\T^2} u \DeltaH u 
		&= \int_{\T^2} \nablaH  u^\perp \nablaH u
		= \int_{\T^2} (\nablaH  u)^\perp \nablaH u
		= 0 . 
	\end{align*}
	By absorbing this implies 
	\begin{equation*}
		\| \partial_t u \|_{\rL^2(\T^2)}^2 
		+
		\sqrt{\nu_{\rH}}\cdot \partial_t \| \nablaH u \|_{\rL^2(\T^2)}^2
		+ \nu_{\rH}\cdot \|\DeltaH u \|_{\rL^2(\T^2)}^2
		\leq 
		C \cdot \| k \|_{\rL^2(\T^2)}^2
		, 
	\end{equation*}
	The pressure satisfies the elliptic equation
	\begin{equation*}
		\DeltaH \pi = \divH k - \divH u^\perp  
	\end{equation*}
	which yields to the estimate
	\begin{align*}
		\| \nablaH \pi \|_{\rL^2(\Omega)}
		&\leq C \cdot (\| \divH k \|_{\rH^{-1}(\T^2)}
		+ \| \divH u^\perp \|_{\rH^{-1}(\T^2)}) \\
		&\leq C \cdot (\| k \|_{\rL^2(\T^2)}
		+ \| u \|_{\rL^2(\T^2)}) ,
	\end{align*}
	where we have used the characterization of the negative order Sobolev space
	\begin{equation*}
		\rH^{-1}(\T^2) = \{ \divH k \colon k \in \rL^2(\T^2) \} .
	\end{equation*}
	Adding both estimates yields the claim. 
\end{proof}

\end{document}
